\documentclass[11pt,a4paper]{article}
% Shared preamble for RuPhO English problem statements (match T1 2026 style).
% Use: \input{latex/rupho_en_preamble.tex} after \documentclass.
\usepackage[margin=1in]{geometry}
\usepackage{amsmath,amssymb}
\usepackage{graphicx}
\usepackage{float}
\usepackage{enumitem}
\usepackage{microtype}
\usepackage{caption}
\usepackage{tabularx}
\usepackage{hyperref}

\captionsetup{font=small,labelfont=bf,skip=6pt}

\setlist[itemize]{leftmargin=1.2cm}
\setlist[enumerate]{leftmargin=1.2cm}

% One outer box per subproblem: [id | statement | points] (no inner rules)
\newcommand{\problembox}[3]{%
  \par\addvspace{0.42em}%
  \noindent
  \setlength{\fboxsep}{7.5pt}%
  \setlength{\fboxrule}{0.95pt}%
  \fbox{%
    \begin{minipage}{\dimexpr\linewidth-2\fboxsep-2\fboxrule\relax}
    \begin{tabularx}{\linewidth}{@{}>{\raggedright\arraybackslash\bfseries}p{2.1em} @{\hspace{0.9em}} >{\raggedright\arraybackslash}X @{\hspace{0.9em}} >{\raggedleft\arraybackslash\bfseries}p{2.4em}@{}}
    #1 & #2 & #3
    \end{tabularx}
    \end{minipage}%
  }%
  \par\addvspace{0.32em}%
}


\title{\textbf{Oscillations in capacitors}}
\author{\normalsize Y14 (2014) --- English translation}
\date{}
\begin{document}
\maketitle

\section*{Part A. Rapid Oscillations}

A free molecule of hydrochloric acid ($\rm HCl$) is an electric dipole of length $L = 0.0215~\text{nm}$ and with atomic charges $+q$ and $-q$ ($q = 1.6 \cdot 10^{-19}~\text{C}$). The molecule was placed in a uniform electric field of strength $E = 30~\text{kV}/\text{cm}$. The masses of hydrogen and chlorine atoms are $m_1 = 1.67 \cdot 10^{-27}~\text{g}$ and $m_2 = 59.2 \cdot 10^{-27}~\text{g}$, respectively.

\problembox{A1}{Find the period $T_0$ of small oscillations of the dipole.}{2.00}

\problembox{A2}{Estimate (with an accuracy no worse than $5\,\%$) the minimum time $\tau$ during which the dipole will rotate from a position in which its axis is perpendicular to the electric field lines to a position in which this axis is parallel to the field lines.}{2.00}

Note: When the dipole deviates through a large angle $\varphi_\text{max}$, the period of its oscillations will depend on the angle $\varphi_\text{max}$ according to the law: \[ T = T_0 \left[ 1 + \sum\limits_{k=1}^{\infty} \left( \frac{(2k-1)!!}{(2k)!!} \right)^2 \sin^{2k} \left( \frac{\varphi_\text{max}}{2} \right) \right] \] Double factorial is defined as follows: $k!!=k\cdot(k-2)\cdot \dots \cdot a$, where $a=2$ for even $k$ and $a=1$ for odd $k$.

\[ T = T_0 \left[ 1 + \sum\limits_{k=1}^{\infty} \left( \frac{(2k-1)!!}{(2k)!!} \right)^2 \sin^{2k} \left( \frac{\varphi_\text{max}}{2} \right) \right] \]

The double factorial is defined as: $k!!=k\cdot(k-2)\cdot \dots \cdot a$, where $a=2$ for even $k$ and $a=1$ for odd $k$.

\section*{Part B. Slow Oscillations}

In the chamber, evacuated to a high vacuum, there is a flat capacitor whose plates are vertical. Distance between plates $D$. In the center of the capacitor there is a small ball (its specific charge is $q_1$). An alternating voltage $U=U_0 \sin \omega t$ is applied to the capacitor plates. At the moment when the voltage went to zero, the ball was released. After some time $\tau$ the ball collided with the capacitor plate. The experiment was repeated, but this time a voltage $U = U_0 \sin^2 \omega t$ was applied to the capacitor. At the moment when the voltage went to zero, the ball was released again and it again hit the same place on the capacitor plate.

\begin{figure}[H]
    \centering
    \includegraphics[width=0.72\textwidth]{T10_EN_assets/fig1.jpg}
    \caption{Figure 1.}
\end{figure}

\problembox{B1}{Determine the frequency $\omega$ of the AC voltage source if $T \gg \tau$ is known. Here $T$ is the period of voltage oscillations at the generator.}{10.00}

\vspace{1em}
\noindent\textit{Source:} \url{https://pho.rs/p/3990}

\end{document}